% TOP_PROBLEM: {"id": 30003534, "title": "$L^p$ Bounds for the Triangular Hilbert Transform", "category_id": 8, "category": {"id": 8, "name": "analysis", "display_name": "Analysis", "description": "Limits, continuity, calculus, and function theory.", "slug": "analysis", "order_index": 8, "created_at": "2026-07-31T15:26:25.671Z"}, "status": "open"}
% FIRST_POSTED: 2026-10-06
% ATTEMPT_STATUS: unresolved
% OPENAI_CLAIM: {"schema_version":1,"source_repo":"https://github.com/openai/math","source_commit":"adc7f1241b42e322a6451854ab7e4b4c146bf78a","release_date":"2026-10-06","match":"partial","resolution":"partial","independently_reviewed":false,"second_pass":"n/a","families":[{"family":"082","title":"Annular variation and dyadic absolute bounds for the triangular Hilbert transform","subject":"Real and complex analysis","manuscripts":[{"title":"An L³ bound for the dyadic triangular Hilbert form","pdf_path":"preprints/An-L3-bound-for-the-dyadic-triangular-Hilbert-form-October-5-2026/dyadic-triangular-hilbert.pdf","readme_path":"preprints/An-L3-bound-for-the-dyadic-triangular-Hilbert-form-October-5-2026/README.md","date":"2026-10-05","theorem_ref":"thm:main (sections/01-introduction.tex:39)","statement_sources":[{"path":"preprints/An-L3-bound-for-the-dyadic-triangular-Hilbert-form-October-5-2026/build/sections/01-introduction.tex","line":39,"label":"thm:main"}]},{"title":"Annular variation of the triangular Hilbert transform at the symmetric point","pdf_path":"preprints/Annular-variation-of-the-triangular-Hilbert-transform-at-the-symmetric-point-October-5-2026/annular-variation.pdf","readme_path":"preprints/Annular-variation-of-the-triangular-Hilbert-transform-at-the-symmetric-point-October-5-2026/README.md","date":"2026-10-05","theorem_ref":"thm:variation (sections/introduction.tex:27)","statement_sources":[{"path":"preprints/Annular-variation-of-the-triangular-Hilbert-transform-at-the-symmetric-point-October-5-2026/build/sections/introduction.tex","line":27,"label":"thm:variation"}]},{"title":"The maximal triangular Hilbert transform at the symmetric point","pdf_path":"preprints/The-maximal-triangular-Hilbert-transform-at-the-symmetric-point-September-24-2026/paper.pdf","readme_path":"preprints/The-maximal-triangular-Hilbert-transform-at-the-symmetric-point-September-24-2026/README.md","date":"2026-09-24","theorem_ref":"thm:main (sections/introduction.tex:16)","statement_sources":[{"path":"preprints/The-maximal-triangular-Hilbert-transform-at-the-symmetric-point-September-24-2026/build/sections/introduction.tex","line":16,"label":"thm:main"}]}],"lean":{"doc_path":"lean/docs/082.md","comparators":["lean/ComparatorChallenges/TriangularHilbert.json","lean/ComparatorChallenges/TriangularHilbert.lean"],"declarations":["OAI.TriangularHilbert.main_estimate"],"covers_main_theorem":false},"reasoning_trace":null}]}
\documentclass[11pt]{article}
\usepackage[margin=25mm]{geometry}
\usepackage{fontspec}
\setmainfont{Latin Modern Roman}
\usepackage{amsmath,amssymb,mathtools}
\usepackage{unicode-math}
\setmathfont{Latin Modern Math}
\usepackage{enumitem}
\usepackage{xurl,hyperref,bookmark}
\hypersetup{colorlinks=true,urlcolor=blue}
\setcounter{secnumdepth}{0}
\setlength{\parindent}{0pt}
\setlength{\parskip}{5pt}
\setlength{\emergencystretch}{3em}
\begin{document}
\section*{\$L\textasciicircum{}p\$ Bounds for the Triangular Hilbert Transform}

\subsection{OpenAI's claim}
\textit{Dyadic triangular Hilbert form}: The dyadic manuscript claims a uniform bound for the sum of the absolute local triangular contributions over any finite selection of scales; the three inputs are arbitrary real $L^3$ functions, and local coefficients may vary independently with modulus at most one.

\textit{Annular triangular variation}: The paper announces an $r$-variation estimate for every $r>2$ from two complex $L^3$ inputs to $L^{3/2}$; output-dependent partitions over positive annular scales are allowed, and the result implies maximal and principal-value conclusions.

\textit{Maximal triangular Hilbert transform}: The announced symmetric triangular Hilbert estimate controls the supremum over both finite hard cutoffs for arbitrary complex $L^3$ inputs; it also gives joint principal-value convergence almost everywhere and in $L^{3/2}$.

\subsection{Mathematical statement of the claimed result}
\paragraph{Mathematical claim 1.} For $k\in\mathbb Z$, let $\mathcal D_k=\{[2^km,2^k(m+1)):m\ge0\}$, let $h_I$ equal $1$ on the left half of $I$ and $-1$ on its right half, and let $\mathcal A_k$ comprise triples whose interval indices have bitwise XOR zero. Define \[L_{\mathbf I}=2^{-k}\int_{I_0\times I_1\times I_2}F_0(x_0,x_1)F_1(x_1,x_2)F_2(x_2,x_0)\prod_{v=0}^2h_{I_v}(x_v)\,dx_0\,dx_1\,dx_2.\] For real bounded compactly supported measurable $F_v$ on $[0,\infty)^2$ and any finite $S\subset\mathbb Z$, \[\sum_{k\in S}\sum_{\mathbf I\in\mathcal A_k}|L_{\mathbf I}|\le40\prod_{v=0}^2\|F_v\|_3.\] This is the dyadic XOR model, not the continuous triangular transform.

\paragraph{Mathematical claim 2.} For complex $F,G\in L^3(\mathbb R^2)$ set $B_{a,b}(F,G)(x,y)=\int_{a<|t|<b}F(x+t,y)G(x,y+t)\,dt/t$. For $r>2$ define \[V_r(F,G)=\sup_{J\ge1,\ 0<t_0<\cdots<t_J}\left(\sum_{j=1}^J|B_{t_{j-1},t_j}(F,G)|^r\right)^{1/r}.\] The supremum may be taken over rational endpoints and imposes no bound on $J$. The claim is $\|V_r(F,G)\|_{3/2}\le C_r\|F\|_3\|G\|_3$, with quantities defined off one common null set.

\paragraph{Mathematical claim 3.} For $F,G\in L^3(\mathbb R^2;\mathbb C)$ let \[B_*(F,G)(x,y)=\sup_{0<a<b<\infty}\left|\int_{a<|t|<b}F(x+t,y)G(x,y+t)\,\frac{dt}{t}\right|.\] There is an absolute finite $C$ with $\|B_*(F,G)\|_{L^{3/2}(\mathbb R^2)}\le C\|F\|_3\|G\|_3$. Both finite hard-truncation endpoints vary in the pointwise supremum; this is the balanced exponent pair $(3,3)$.

\subsection{Outline of the argument}
\textit{Dyadic triangular Hilbert form}: Fine dyadic step functions reduce the inputs to matrices between atom spaces. At each of the three variables, the two incident matrices are joined into one row operator and assigned a positive trace energy. A local bound charges the form to the energy increment between a triple and its children; telescoping and endpoint norm bounds control all selected scales.

\textit{Annular triangular variation}: The proof first bounds the sum of absolute contributions from finitely many disjoint annuli, selected separately at each output location. Quantitative growth control for this counting quantity yields the variation inequality. Approximation and the maximal consequence then give common almost-everywhere limits and convergence in the output norm.

\textit{Maximal triangular Hilbert transform}: The cyclic form is represented by matrix traces and a positive energy evolving under Gaussian smoothing. Mixed trace estimates compare energy derivatives at adjacent variables. A further bound pays for entries switched on or off at separately chosen scales, so telescoping remains uniform under pointwise selection of the cutoffs.

\subsection{Scope relative to this problem}
\textit{Dyadic triangular Hilbert form}: The paper proves one symmetric L\textasciicircum{}3 triangular bound (annular variation/maximal continuous transform, or its dyadic form). The question asks the full range of exponent triples for the continuous triangular Hilbert transform; the dyadic model also does not automatically transfer to all continuous kernels.

\textit{Annular triangular variation}: The paper proves one symmetric L\textasciicircum{}3 triangular bound (annular variation/maximal continuous transform, or its dyadic form). The question asks the full range of exponent triples for the continuous triangular Hilbert transform; the dyadic model also does not automatically transfer to all continuous kernels.

\textit{Maximal triangular Hilbert transform}: The paper proves one symmetric L\textasciicircum{}3 triangular bound (annular variation/maximal continuous transform, or its dyadic form). The question asks the full range of exponent triples for the continuous triangular Hilbert transform; the dyadic model also does not automatically transfer to all continuous kernels.

This page records partial support. Complete resolution of the catalogue question is not established by this record.

\subsection{Formal verification}
Family 082 describes Lean coverage for the two-cutoff symmetric maximal estimate for complex inputs; later variation and dyadic absolute-sum refinements are separate.

\textit{Dyadic triangular Hilbert form}: No manuscript-specific Comparator is listed in the pinned inventory.

The documented selection does not cover the entire manuscript assertion relevant here.

\textit{Annular triangular variation}: No manuscript-specific Comparator is listed in the pinned inventory.

The documented selection does not cover the entire manuscript assertion relevant here.

\textit{Maximal triangular Hilbert transform}: The manuscript links \texttt{TriangularHilbert.lean} (declaration \texttt{OAI.TriangularHilbert.main\_estimate}).

The documented selection includes this manuscript’s main assertion; that assertion retains the restrictions stated above.

This site has not run the Lean development. These links describe OpenAI’s selected formal statements and should be read with their stated scope.

\subsection{Status caveat}
OpenAI reports that an internal model produced the manuscripts. Their arguments have not been independently verified as proofs by this site. Any portion outside the documented formal scope may contain mistakes, and the catalogue label reports the announced mathematical claim.

\subsection{Sources}
\begin{itemize}
\item OpenAI, \emph{An L\({}^{3}\) bound for the dyadic triangular Hilbert form}, 2026-10-05; citation key \texttt{OAI:An-L3-bound-for-the-dyadic-triangular-Hilbert-form-October-5-2026}. \href{https://github.com/openai/math/blob/adc7f1241b42e322a6451854ab7e4b4c146bf78a/preprints/An-L3-bound-for-the-dyadic-triangular-Hilbert-form-October-5-2026/dyadic-triangular-hilbert.pdf}{PDF (GitHub)}; \href{https://github.com/openai/math/raw/adc7f1241b42e322a6451854ab7e4b4c146bf78a/preprints/An-L3-bound-for-the-dyadic-triangular-Hilbert-form-October-5-2026/dyadic-triangular-hilbert.pdf}{PDF (direct)}; \href{https://github.com/openai/math/blob/adc7f1241b42e322a6451854ab7e4b4c146bf78a/preprints/An-L3-bound-for-the-dyadic-triangular-Hilbert-form-October-5-2026/README.md}{Manuscript page and citation}; \href{https://github.com/openai/math/blob/main/preprints/An-L3-bound-for-the-dyadic-triangular-Hilbert-form-October-5-2026/dyadic-triangular-hilbert.pdf}{Latest version}; Statement source: \href{https://github.com/openai/math/blob/adc7f1241b42e322a6451854ab7e4b4c146bf78a/preprints/An-L3-bound-for-the-dyadic-triangular-Hilbert-form-October-5-2026/build/sections/01-introduction.tex\#L39}{thm:main}.
\item OpenAI, \emph{Annular variation of the triangular Hilbert transform at the symmetric point}, 2026-10-05; citation key \texttt{OAI:Annular-variation-of-the-triangular-Hilbert-transform-at-the-symmetric-point-October-5-2026}. \href{https://github.com/openai/math/blob/adc7f1241b42e322a6451854ab7e4b4c146bf78a/preprints/Annular-variation-of-the-triangular-Hilbert-transform-at-the-symmetric-point-October-5-2026/annular-variation.pdf}{PDF (GitHub)}; \href{https://github.com/openai/math/raw/adc7f1241b42e322a6451854ab7e4b4c146bf78a/preprints/Annular-variation-of-the-triangular-Hilbert-transform-at-the-symmetric-point-October-5-2026/annular-variation.pdf}{PDF (direct)}; \href{https://github.com/openai/math/blob/adc7f1241b42e322a6451854ab7e4b4c146bf78a/preprints/Annular-variation-of-the-triangular-Hilbert-transform-at-the-symmetric-point-October-5-2026/README.md}{Manuscript page and citation}; \href{https://github.com/openai/math/blob/main/preprints/Annular-variation-of-the-triangular-Hilbert-transform-at-the-symmetric-point-October-5-2026/annular-variation.pdf}{Latest version}; Statement source: \href{https://github.com/openai/math/blob/adc7f1241b42e322a6451854ab7e4b4c146bf78a/preprints/Annular-variation-of-the-triangular-Hilbert-transform-at-the-symmetric-point-October-5-2026/build/sections/introduction.tex\#L27}{thm:variation}.
\item OpenAI, \emph{The maximal triangular Hilbert transform at the symmetric point}, 2026-09-24; citation key \texttt{OAI:The-maximal-triangular-Hilbert-transform-at-the-symmetric-point-September-24-2026}. \href{https://github.com/openai/math/blob/adc7f1241b42e322a6451854ab7e4b4c146bf78a/preprints/The-maximal-triangular-Hilbert-transform-at-the-symmetric-point-September-24-2026/paper.pdf}{PDF (GitHub)}; \href{https://github.com/openai/math/raw/adc7f1241b42e322a6451854ab7e4b4c146bf78a/preprints/The-maximal-triangular-Hilbert-transform-at-the-symmetric-point-September-24-2026/paper.pdf}{PDF (direct)}; \href{https://github.com/openai/math/blob/adc7f1241b42e322a6451854ab7e4b4c146bf78a/preprints/The-maximal-triangular-Hilbert-transform-at-the-symmetric-point-September-24-2026/README.md}{Manuscript page and citation}; \href{https://github.com/openai/math/blob/main/preprints/The-maximal-triangular-Hilbert-transform-at-the-symmetric-point-September-24-2026/paper.pdf}{Latest version}; Statement source: \href{https://github.com/openai/math/blob/adc7f1241b42e322a6451854ab7e4b4c146bf78a/preprints/The-maximal-triangular-Hilbert-transform-at-the-symmetric-point-September-24-2026/build/sections/introduction.tex\#L16}{thm:main}.
\item Family 082: \href{https://github.com/openai/math/blob/adc7f1241b42e322a6451854ab7e4b4c146bf78a/lean/docs/082.md}{Lean scope}.
\item \href{https://github.com/openai/math/blob/adc7f1241b42e322a6451854ab7e4b4c146bf78a/lean/ComparatorChallenges/TriangularHilbert.json}{Comparator TriangularHilbert.json}.
\item \href{https://github.com/openai/math/blob/adc7f1241b42e322a6451854ab7e4b4c146bf78a/lean/ComparatorChallenges/TriangularHilbert.lean}{Comparator TriangularHilbert.lean}.
\item \href{https://github.com/openai/math/blob/adc7f1241b42e322a6451854ab7e4b4c146bf78a/CONTENTS.md}{OpenAI catalogue entry}.
\end{itemize}
\end{document}
